% Part: incompleteness
% Chapter: arithmetization-syntax
% Section: proofs-in-lk

\documentclass[../../../include/open-logic-section]{subfiles}

\begin{document}

\olfileid[id]{inc}{art}{plk}
\olsection{\usetoken{P}{derivation} dalam $\Log{LK}$}

\begin{explain}
Untuk mengaritmetisasi !!{derivation}s, kita harus merepresentasikan
!!{derivation}s sebagai bilangan. Karena !!{derivation}s merupakan pohon sekuen
dan setiap inferensi juga memiliki label, pendekatan yang paling jelas ialah
representasi rekursif: kita merepresentasikan !!{derivation} sebagai tupel yang
komponen-komponennya adalah sekuen akhir, label, dan representasi
sub-!!{derivation}s yang menuju
premis-premis inferensi terakhir.
\end{explain}

\begin{defn}
Jika $\Gamma$ adalah barisan berhingga !!{sentence}s,
$\Gamma = \tuple{!A_1, \dots, !A_n}$, maka
$\Gn{\Gamma} = \tuple{\Gn{!A_1}, \dots, \Gn{!A_n}}$.

Jika $\Gamma \Sequent \Delta$ adalah sekuen, maka bilangan G\"odel dari
$\Gamma \Sequent \Delta$ ialah
\[
\Gn{\Gamma \Sequent \Delta} = \tuple{\Gn{\Gamma}, \Gn{\Delta}}
\]

Jika $\pi$ adalah !!a{derivation} dalam $\Log{LK}$, maka $\Gn{\pi}$
didefinisikan sebagai berikut:
\begin{enumerate}
\item Jika $\pi$ hanya terdiri atas sekuen awal $\Gamma \Sequent \Delta$,
  maka $\Gn{\pi}$ ialah
  \[
    \tuple{0, \Gn{\Gamma \Sequent \Delta}}.
  \]
\item Jika $\pi$ berakhir dengan inferensi yang memiliki satu atau dua premis,
  sekuen kesimpulannya ialah $\Gamma \Sequent \Delta$, dan $\pi_1$ serta
  $\pi_2$ adalah subpembuktian langsung yang berakhir dengan premis inferensi
  terakhir, maka $\Gn{\pi}$ masing-masing ialah
  \begin{align*}
    & \tuple{1, \Gn{\pi_1}, \Gn{\Gamma \Sequent \Delta}, k} \text{ atau}\\
    & \tuple{2, \Gn{\pi_1}, \Gn{\pi_2}, \Gn{\Gamma \Sequent \Delta}, k},
  \end{align*}
  dengan $k$ diberikan oleh tabel berikut menurut aturan yang digunakan dalam
  inferensi terakhir:

  \begin{tabular}{lccccccc}
    \text{Aturan:} & \LeftR{\Weakening} & \RightR{\Weakening} &
    \LeftR{\Contraction} & \RightR{\Contraction} &
    \LeftR{\Exchange} & \RightR{\Exchange} \\
    $k$: & 1 & 2 & 3 & 4 & 5 & 6 \\[2ex]
    \text{Aturan:} & \LeftR{\lnot} & \RightR{\lnot} &
    \LeftR{\land} & \RightR{\land} &
    \LeftR{\lor} & \RightR{\lor} \\
    $k$: & 7 & 8 & 9 & 10 & 11 & 12 \\[2ex]
    \text{Aturan:} & \LeftR{\lif} & \RightR{\lif} &
    \LeftR{\lforall} & \RightR{\lforall} &
    \LeftR{\lexists} & \RightR{\lexists} \\
    $k$: & 13 & 14 & 15 & 16 & 17 & 18 \\[2ex]
    \text{Aturan:} & \Cut & = \\
    $k$: & 19 & 20
  \end{tabular}
\end{enumerate}
\end{defn}

\begin{ex}
  Tinjaulah !!{derivation} yang sangat sederhana berikut:
  \begin{prooftree}
    \Axiom$!A \fCenter !A$
    \RightLabel{\LeftR{\land}}
    \UnaryInf$!A \land !B \fCenter !A$
    \RightLabel{\RightR{\lif}}
    \UnaryInf$\fCenter (!A \land !B) \lif !A$
  \end{prooftree}
  Bilangan G\"odel sekuen awalnya ialah $p_0 = \tuple{0,
    \Gn{!A \Sequent !A}}$. Bilangan G\"odel !!{derivation} yang berakhir
  pada kesimpulan~$\LeftR{\land}$ ialah $p_1 =
  \tuple{1, p_0, \Gn{!A \land !B \Sequent !A}, 9}$ ($1$ karena
  $\LeftR{\land}$ memiliki satu premis, lalu bilangan G\"odel kesimpulan
  $!A \land !B \Sequent !A$, dan $9$ adalah bilangan yang mengodekan
  $\LeftR{\land}$). Jadi, bilangan G\"odel seluruh !!{derivation} itu ialah
  $\tuple{1, p_1, \Gn{\Sequent (!A \land !B) \lif !A)}, 14}$, yakni
  \[
  \tuple{1, \tuple{1, \tuple{0, \Gn{!A \Sequent !A}}, \Gn{!A \land !B \Sequent !A}, 9},
    \Gn{\Sequent (!A \land !B) \lif !A}, 14}.
  \]
\end{ex}

\begin{explain}
Setelah menetapkan suatu representasi bagi !!{derivation}s, kita juga harus
menunjukkan bahwa kita dapat memanipulasi !!{derivation}s tersebut secara
rekursif primitif serta menyatakan sifat dan relasi pentingnya dengan cara yang
sama. Beberapa operasi sederhana: misalnya, jika diberikan bilangan
G\"odel~$p$ dari !!a{derivation}, $\fn{EndSequent}(p) = (p)_{(p)_0+1}$ memberikan
bilangan G\"odel sekuen akhirnya dan $\fn{LastRule}(p) = (p)_{(p)_0+2}$
memberikan kode aturan terakhirnya. Sifat $\fn{Sequent}(s)$ yang didefinisikan
oleh
\[
  \len{s} = 2 \land \bforall{i<\len{(s)_0} + \len{(s)_1}}{\fn{Sent}(((s)_0 \concat (s)_1)_i)}
\]
berlaku bagi $s$ jika dan hanya jika $s$ adalah bilangan G\"odel sekuen yang
terdiri atas !!{sentence}s. Sebagian sifat lain jauh lebih sulit. Setidaknya
kita akan membuat sketsa cara menanganinya. Tujuan kita ialah menunjukkan
bahwa relasi ``$\pi$ adalah !!a{derivation} dari~$!A$ berdasarkan~$\Gamma$''
merupakan relasi rekursif primitif antara bilangan G\"odel $\pi$ dan~$!A$.
\end{explain}

\begin{prop}\ollabel{prop:followsby}
  Sifat $\fn{Correct}(p)$, yang berlaku jika dan hanya jika inferensi terakhir
  dalam !!{derivation}~$\pi$ dengan bilangan G\"odel~$p$ benar, bersifat
  rekursif primitif.
\end{prop}

\begin{proof}
  $\Gamma \Sequent \Delta$ adalah sekuen awal jika terdapat
  !!a{sentence}~$!A$ sedemikian rupa sehingga $\Gamma \Sequent \Delta$ adalah
  $!A \Sequent !A$, atau terdapat suku~$t$ sedemikian rupa sehingga
  $\Gamma \Sequent \Delta$ adalah $\emptyset \Sequent \eq[t][t]$. Dalam
  bilangan G\"odel, $\fn{InitSeq}(s)$ berlaku jika dan hanya jika
  \begin{align*}
  \bexists{x < s}{} (\fn{Sent}(x) & \land
  s = \tuple{\tuple{x},\tuple{x}})
  \lor {}\\
  \bexists{t<s}{} (\fn{Term}(t) & \land
  s = \tuple{0, \tuple{\Gn{{\eq}(} \concat t \concat \Gn{,} \concat t \concat \Gn{)}}}).
  \end{align*}

  Kita juga harus menunjukkan bahwa bagi setiap aturan inferensi~$R$, relasi
  $\fn{FollowsBy}_R(p)$ bersifat rekursif primitif, dengan
  $\fn{FollowsBy}_R(p)$ berlaku jika dan hanya jika $p$ adalah bilangan G\"odel
  !!{derivation}~$\pi$, dan sekuen akhir~$\pi$ mengikuti melalui penerapan~$R$
  yang benar dari sub-!!{derivation}s langsung~$\pi$.

  Salah satu kasus sederhana ialah aturan \RightR{\land}. Jika $\pi$
  berakhir dengan inferensi $\RightR{\land}$ yang benar, bentuknya ialah:
  \begin{prooftree}
    \AxiomC{}
    \RightLabel{$\pi_1$}
    \Deduce$\Gamma \fCenter \Delta, !A$

    \AxiomC{}
    \RightLabel{$\pi_2$}
    \Deduce$\Gamma \fCenter \Delta, !B$

    \RightLabel{\RightR\land}
    \BinaryInf$\Gamma \fCenter \Delta, !A \land !B$
  \end{prooftree}
  Jadi, inferensi terakhir dalam !!{derivation} $\pi$ merupakan penerapan
  $\RightR{\land}$ yang benar jika dan hanya jika terdapat barisan
  !!{sentence}s $\Gamma$ dan $\Delta$, serta dua !!{sentence}s $!A$ dan~$!B$,
  sedemikian rupa sehingga sekuen akhir $\pi_1$ ialah
  $\Gamma \Sequent \Delta, !A$, sekuen akhir $\pi_2$ ialah
  $\Gamma \Sequent \Delta, !B$, dan sekuen akhir $\pi$ ialah
  $\Gamma \Sequent \Delta, !A \land !B$. Kita hanya perlu menerjemahkannya ke
  dalam bilangan G\"odel. Jika $s = \Gn{\Gamma \Sequent \Delta}$, maka
  $(s)_0 = \Gn{\Gamma}$ dan $(s)_1 = \Gn{\Delta}$. Jadi,
  $\fn{FollowsBy}_{\RightR{\land}}(p)$ berlaku jika dan hanya jika
\begin{align*}
& \bexists{g < p}{\bexists{d < p}{\bexists{a < p}{\bexists{b < p}{\quad}}}} \\
& \qquad \fn{EndSequent}(p) =
  \tuple{g, d \concat
    \tuple{\Gn{(} \concat a \concat \Gn{\land} \concat b \concat \Gn{)}}} \land {} \\
& \qquad \fn{EndSequent}((p)_1) = \tuple{g, d \concat \tuple{a}} \land {}\\
& \qquad \fn{EndSequent}((p)_2) = \tuple{g, d \concat \tuple{b}} \land {}\\
& \qquad (p)_0 = 2 \land \fn{LastRule}(p) = 10.
\end{align*}
Baris-baris tersebut masing-masing menyatakan bahwa ``terdapat barisan
($\Gamma$) dengan bilangan G\"odel~$g$, terdapat barisan ($\Delta$) dengan
bilangan G\"odel~$d$, !!a{formula} ($!A$) dengan bilangan G\"odel~$a$, dan
!!a{formula} ($!B$) dengan bilangan G\"odel~$b$,'' sedemikian rupa sehingga
``sekuen akhir $\pi$ ialah $\Gamma \Sequent \Delta, !A \land !B$,'' ``sekuen
akhir $\pi_1$ ialah $\Gamma \Sequent \Delta, !A$,'' ``sekuen akhir $\pi_2$
ialah $\Gamma \Sequent \Delta, !B$,'' dan ``$\pi$ memiliki dua subderivasi
langsung dan aturan inferensi terakhirnya ialah $\RightR\land$ (dengan
bilangan~$10$).''

Inferensi terakhir dalam~$\pi$ merupakan penerapan $\RightR\lexists$ yang
benar jika dan hanya jika terdapat barisan $\Gamma$ dan $\Delta$,
!!a{formula}~$!A$, variabel~$x$, dan suku~$t$, sedemikian rupa sehingga sekuen
akhir $\pi$ ialah $\Gamma \Sequent \Delta, \lexists[x][!A]$ dan sekuen akhir
$\pi_1$ ialah $\Gamma \Sequent \Delta, \Subst{!A}{t}{x}$. Jadi, dalam bilangan
G\"odel, $\fn{FollowsBy}_{\RightR\lexists}(p)$ berlaku jika dan hanya jika
\begin{align*}
  & \bexists{g<p}{\bexists{d<p}{\bexists{a<p}{\bexists{x<p}{\bexists{t<p}{\quad}}}}}\\
  & \qquad \fn{EndSequent}(p) =
  \tuple{
    g,
    d \concat \tuple{\Gn{\lexists} \concat x \concat a}
  } \land {}\\
  & \qquad \fn{EndSequent}((p)_1) =
  \tuple{
    g,
    d \concat \tuple{\fn{Subst}(a, t, x)}
  } \land {}\\
  & \qquad (p)_0 = 1 \land \fn{LastRule}(p) = 18.
\end{align*}
Kemudian kita mendefinisikan $\fn{Correct}(p)$ sebagai
\begin{multline*}
  \fn{Sequent}(\fn{EndSequent}(p)) \land {}\\
  [(\fn{LastRule}(p) = 1 \land
  \fn{FollowsBy}_{\LeftR\Weakening}(p)) \lor \dots \lor {}\\
  (\fn{LastRule}(p) = 20 \land \fn{FollowsBy}_{\eq}(p)) \lor {}\\
  (p)_0 = 0 \land \fn{InitSeq}(\fn{EndSequent}(p))]
\end{multline*}
Baris pertama memastikan bahwa sekuen akhir derivasi berkode~$p$ benar-benar merupakan sekuen
yang terdiri atas !!{sentence}s. Baris terakhir mencakup kasus ketika $p$
hanyalah sekuen awal.
\end{proof}

\begin{prob}
  Definisikan sifat-sifat berikut seperti dalam
  \olref[inc][art][plk]{prop:followsby}:
  \begin{enumerate}
  \item $\fn{FollowsBy}_{\Cut}(p)$;
  \item $\fn{FollowsBy}_{\LeftR\lif}(p)$;
  \item $\fn{FollowsBy}_{\eq}(p)$;
  \item $\fn{FollowsBy}_{\RightR{\lforall}}(p)$.
  \end{enumerate}
  Untuk yang terakhir, Anda juga harus menunjukkan bahwa kita dapat menguji
  secara rekursif primitif apakah inferensi terakhir dari !!{derivation}
  dengan bilangan G\"odel~$p$ memenuhi syarat variabel eigen, yakni variabel
  eigen~$a$ dari $\RightR{\lforall}$ tidak muncul dalam sekuen akhir.
\end{prob}

\begin{prop}
  \ollabel{prop:deriv}
  Relasi $\fn{Deriv}(p)$, yang berlaku jika $p$ adalah bilangan G\"odel suatu
  !!{derivation}~$\pi$ yang benar, bersifat rekursif primitif.
\end{prop}

\begin{proof}
  !!^a{derivation}~$\pi$ benar jika setiap inferensinya merupakan penerapan
  aturan yang benar, yakni jika setiap sub-!!{derivation}s berakhir dengan
  inferensi yang benar. Jadi, $\fn{Deriv}(p)$ jika dan hanya jika
  \[
  \bforall{i<\len{\fn{SubtreeSeq}(p)}}{\fn{Correct}((\fn{SubtreeSeq}(p))_i)}.
  \]
\end{proof}


\begin{prop}
Andaikan $\Gamma$ adalah himpunan rekursif primitif yang terdiri atas
!!{sentence}s. Maka relasi $\Prf[\Gamma](x, y)$ yang menyatakan ``$x$ adalah
kode !!a{derivation}~$\pi$ dari $\Gamma_0 \Sequent !A$ untuk suatu
$\Gamma_0 \subseteq \Gamma$ yang berhingga, dan $y$ adalah bilangan G\"odel
dari~$!A$'' bersifat rekursif primitif.
\end{prop}

\begin{proof}
Andaikan ``$y \in \Gamma$'' diberikan oleh predikat rekursif primitif
$R_\Gamma(y)$. Kita harus menunjukkan bahwa $\Prf[\Gamma](x, y)$ bersifat
rekursif primitif, dengan relasi ini berlaku jika dan hanya jika $y$ adalah
bilangan G\"odel suatu kalimat~$!A$ dan $x$ adalah kode suatu
$\Log{LK}$-!!{derivation} dengan sekuen akhir $\Gamma_0 \Sequent !A$.

Menurut proposisi sebelumnya, sifat $\fn{Deriv}(x)$, yang berlaku jika dan
hanya jika $x$ adalah kode !!{derivation}~$\pi$ yang benar dalam $\Log{LK}$,
bersifat rekursif primitif. Jika $x$ adalah kode semacam itu, maka
$\fn{EndSequent}(x)$ adalah kode sekuen akhir~$\pi$, sehingga
$(\fn{EndSequent}(x))_0$ adalah kode sisi kiri sekuen akhir dan
$(\fn{EndSequent}(x))_1$ adalah kode sisi kanannya. Jadi, kita dapat menyatakan
``sisi kanan sekuen akhir~$\pi$ ialah~$!A$'' sebagai
$\len{(\fn{EndSequent}(x))_1} = 1 \land ((\fn{EndSequent}(x))_1)_0 = y$.
Sisi kiri sekuen akhir $\pi$ tentu secara otomatis berhingga; kita hanya perlu
menyatakan bahwa setiap kalimat di dalamnya termasuk dalam~$\Gamma$. Dengan
demikian, kita dapat mendefinisikan $\Prf[\Gamma](x, y)$ melalui
\begin{align*}
\Prf[\Gamma](x, y) \defiff {}&
\fn{Deriv}(x) \land {} \\
& \bforall{i <
  \len{(\fn{EndSequent}(x))_0}}{R_\Gamma(((\fn{EndSequent}(x))_0)_i)} \land {}\\
& \len{(\fn{EndSequent}(x))_1} = 1 \land ((\fn{EndSequent}(x))_1)_0 = y.
\end{align*}
\end{proof}

\end{document}
